\documentclass[11pt]{article}
\usepackage[margin=1in]{geometry}
\usepackage{amsmath,amssymb,amsthm}
\usepackage{fontspec}
\usepackage[hidelinks,pdftitle={Instance-optimal identity testing at the specified total variation distance},pdfauthor={ProofRush. Research directed by Yoshiaki Otsuka}]{hyperref}
\newtheorem{theorem}{Theorem}[section]
\newtheorem{lemma}[theorem]{Lemma}
\newtheorem{proposition}[theorem]{Proposition}
\newtheorem{corollary}[theorem]{Corollary}
\newcommand{\TV}{\operatorname{TV}}
\newcommand{\E}{\mathbb E}
\newcommand{\Prb}{\mathbb P}
\newcommand{\N}{\mathbb N}
\newcommand{\R}{\mathbb R}
\newcommand{\given}{\,\middle|\,}
\newfontfamily\namefont[Path=fonts/]{ProofRushDirectorJP.otf}
\newcommand{\titleDirector}{{\namefont 大束良明} (Yoshiaki Otsuka)}
\title{Instance-optimal identity testing at the specified total variation distance}
\author{ProofRush\\[2pt]\small AI-assisted research directed by \titleDirector}
\date{27 November 2026\qquad Preprint, version 1}
\begin{document}
\maketitle
\begin{abstract}
Given a known discrete distribution $q$ and $0<e\le1$, let $N_q(e)$ be the least number of independent observations that distinguish $q$ from every distribution at total variation distance strictly greater than $e$, with both errors at most one third. We give a functional of $q$ and $e$ comparable to $N_q(e)$ by absolute constants at the very same $e$. The head consists of the fewest largest atoms whose complement has mass at most $e$. The head contributes a minimized truncated $2/3$-moment. The tail contributes the attainable overlap of a random empirical support at a prescribed testing budget. A second regime near distance one is governed by a clipped $3/2$-moment. The lower bound uses one random channel on a reduced experiment, with the same pilot lists under both hypotheses. A deterministic scale comparison converts the resulting inequalities to a sample-size bound without rescaling the distance. For an explicit finite rational law we give a closed formula using polynomially many bit operations. A finite projection argument covers all discrete laws on a countable alphabet. The source package contains Lean proofs of the two-sided sample-complexity statements.
\end{abstract}

\section{Problem and prior bounds}
Fix a discrete alphabet $\mathcal X$. Its labels form part of the experiment, including any with zero mass under the null. For a specified $q$, a randomized test on $n$ observations is a function $f:\mathcal X^n\to[0,1]$ giving the conditional rejection probability. We seek the least nonnegative integer $N_q(e)$ such that
\[
 \E_{q^{\otimes n}} f\le 1/3,\qquad
 \E_{p^{\otimes n}} f\ge 2/3\quad
 \text{for every probability law }p\text{ with }\TV(p,q)>e .
\]
Here $\TV(p,q)=\tfrac12\sum_i|p_i-q_i|$. Access to the unknown law consists of ordinary independent observations.

Valiant and Valiant \cite{valiantvaliant2017} identified a truncated $2/3$-quasinorm that captures much of the instance dependence. Their lower and upper bounds truncate different multiples of the distance. Blais, Canonne and Gur \cite{BCG} obtained bounds using an inverse Peetre functional and described a harmonic family on which changing the truncation distance changes the sample scale by a polynomial in the alphabet size. Chhor and Carpentier \cite{chhorcarpentier2022} characterized a local minimax separation radius to within constants. Inverting an interval of separation radii can leave an unbounded discrepancy between sample sizes at a specified distance. Canonne formulated the sample-size question explicitly as COLT 2024 Open Question~1 \cite{Canonne}. Dewasurendra \cite{Zipf} subsequently reported same-distance bounds for a Zipfian class under envelope assumptions. Section~6 of that manuscript already discusses relating a weighted head test to sparse changes in the tail. Here a common ordered-record channel supplies both the head and marker obstructions, and a sample-scale dilation closes the comparison for arbitrary nulls. Our comparison constants are uniform over the null distribution and the distance. The use of weighted quadratic tests, adjacent signed perturbations and second-moment arguments builds on the distribution-testing literature. The constructions below connect these methods to the amount of overlap a tail can retain at the original distance.

\section{A functional of the known distribution}\label{sec:functional}
First consider a finite alphabet. If $e\ge 1-\min_i q_i$, return zero. Indeed
$\min(p_i,q_i)\ge p_iq_i$, so $\TV(p,q)=1-\sum_i\min(p_i,q_i)\le1-\min_i q_i$. A point mass at a least-probable label attains equality. The remaining case has $\delta=1-e>0$ and admits a strict alternative. Define
\[
  J=J_\delta(q)=\sum_i q_i\min\{1,\sqrt{q_i/\delta}\}.
\]
If $J\ge16\delta$, let $\Phi(q,e)=\lceil1/J\rceil$.

Otherwise order the probabilities decreasingly and let $H$ be the shortest prefix of mass at least $\delta$. Write $T=\mathcal X\setminus H$, $w=q(T)$, $A=1-w$, $a=\min_{i\in H}q_i$ and $\Delta=e-w$. The last head atom is positive, $0\le\Delta<a$, and $q_i\le a$ for $i\in T$. The largest head label is numbered $1$. For a real $n>0$ define
\begin{align}
 P(n)&=\min_{1\le j\le |H|}\left\{
   \frac{(\sum_{2\le i\le j}q_i^{2/3})^{3/4}}{\sqrt n}
      +\sum_{j<i\le |H|}q_i\right\},\label{eq:P}\\
 K(n,x)&=\sum_{i\in T}\min\{q_i/x,\ n q_i^{3/2}/\sqrt x\}\quad(x>0).\label{eq:K}
\end{align}
If $w=0$, set $X(n)=0$. Otherwise $K(n,x)$ decreases continuously from infinity to zero. Let $X(n)>0$ be its unique unit level. Set
\begin{equation}\label{eq:cross}
 D(n)=\sqrt{wA/n}+1/n+P(n),\qquad
 Y(n)=\Delta+X(n),\qquad c_*=(4\,000\,000)^{-1}.
\end{equation}
Define $\Phi(q,e)$ as the first power of two $n\ge1$ with $D(n)\le c_*Y(n)$. The full off-largest head has no omitted mass, so $D(n)\to0$. Also $Y(1)>0$: either $w<e$, or $w=e>0$ gives a positive root. Since $X$ is nondecreasing, a crossing exists. Equal probabilities may be ordered arbitrarily.

The implicit root has a finite radical expression. Order the tail decreasingly. For every tail prefix $S$, including the empty and full prefix, write $Q_S=\sum_{i\in S}q_i$ and $B_S=n\sum_{i\in T\setminus S}q_i^{3/2}$. Then
\begin{equation}\label{eq:radical}
 \sqrt{X(n)}=\min_S \frac{B_S+\sqrt{B_S^2+4Q_S}}2 .
\end{equation}
For $y=\sqrt{X(n)}>0$, the root equation is $y^2=\sum_{i\in T}\min(q_i,n y q_i^{3/2})$. Splitting this sum at any prefix gives $y^2\le Q_S+B_S y$, with equality for $S=\{i:n y\sqrt{q_i}\ge1\}$. For a zero tail both sides vanish. This calculation also gives an expression using only minimum and maximum up to absolute factors. We use it below to avoid comparisons between sums of radicals on rational inputs.

\begin{theorem}\label{thm:main}
There are universal constants $L,U>0$ such that
\[
   \Phi(q,e)\le L N_q(e),\qquad N_q(e)\le U\Phi(q,e)
\]
for every finite $q$ and $0<e\le1$. One can take $U=49152$. The same constants apply on $\mathbb N$ with $\Phi(q,e)=\sup_m\Phi(q_m,e)$, where $q_m$ is the law of the first $m$ singleton cells and the remaining cell. This supremum is finite. There is an explicit finite rational-input version, computable in time polynomial in the binary description length, with universal comparison constants.
\end{theorem}

For example, take $q=(1-w,w/k,\ldots,w/k)$ with $0<w\le1/2$, $k\ge2$, and choose $e=w$. Then $H$ contains only the largest label and $\Delta=P(n)=0$. Equation~\eqref{eq:K} gives $X(n)=\min(w,n^2w^3/k)$. On writing $t=nw$, one has $D(n)/w=\sqrt{(1-w)/t}+1/t$ and $Y(n)/w=\min(1,t^2/k)$. The crossing occurs at $t\asymp k^{2/5}$, with absolute constants. In this family the total tail mass equals the testing distance exactly.

For each alternative $p$ write $L_H=\sum_{i\in H}(q_i-p_i)_+$ and $O_T=\sum_{i\in T}\min(q_i,p_i)$. The equality
\begin{equation}\label{eq:split-tv}
             \TV(p,q)=w+L_H-O_T
\end{equation}
holds because the total decrease from $q$ to $p$ equals total variation. A strict alternative therefore satisfies $L_H>\Delta+O_T$. This exact difference explains why the tail cannot be discarded at a fixed multiple of $e$.

\section{Upper bounds at the prescribed radius}
\label{sec:upper}

For a set \(T\) of labels, define its overlap and the directed deficit
on its complement by
\[
 O_T(p,q)=\sum_{i\in T}\min(p_i,q_i),\qquad
 L_H(p,q)=\sum_{i\in H}(q_i-p_i)_+,\qquad H=T^\complement .
\]
Writing \(w=q(T)\), normalization gives
\begin{equation}
 \operatorname{TV}(p,q)=w+L_H(p,q)-O_T(p,q).
 \label{eq:upper-head-tail}
\end{equation}
In particular, at \(\Delta=e-w\), the strict alternative is exactly
\(L_H(p,q)>\Delta+O_T(p,q)\).

\subsection{A capped overlap test}

For a positive cap \(b\), replace each selected label \(i\) by
\(m_i=\max(1,\lceil q_i/b\rceil)\) equiprobable copies. Use independent
copy choices under both input laws. On the refined labels,
\(p_i\) and \(q_i\) are divided by the same \(m_i\). The total mass and
overlap on any union of full copy classes are preserved, and every refined
reference probability in the selected set is at most \(b\). Its refined
moment satisfies
\begin{equation}
 \frac{1}{\sqrt{2}}\sum_{i\in T}\min(q_i\sqrt b,q_i^{3/2})
 \ \le\ \sum_{i\in T}\frac{q_i^{3/2}}{\sqrt{m_i}}
 \ \le\ \sum_{i\in T}\min(q_i\sqrt b,q_i^{3/2}).
 \label{eq:upper-clone-moment}
\end{equation}
For \(q_i\le b\) the atom is unchanged. For \(q_i>b\) use
\(q_i/b\le m_i<2q_i/b\). Null-zero labels have one copy and remain
available to the alternative.

We now work on a refined alphabet and denote its probabilities by \(p_i,q_i\).
Suppose \(q_i\le b\) for \(i\in T\), \(0<b\le x\), and
\[
   M=\sum_{i\in T}q_i^{3/2}\ge 8x\sqrt b.
\]
We test against all full laws with \(O_T(p,q)\le x\), with no condition
on the number of tail observations or on \(p(T)\). Put
\(\ell=\lceil8192\sqrt{x}/M\rceil\) and \(r=M^2/x\).
On a first batch of \(\ell\) full observations use the score
\(W(i)=\sqrt{q_i}\) on \(T\) and \(W(i)=0\) off \(T\). Reject if its empirical
mean is below \(7M/8\). Under any law,
\(\operatorname{Var}(W)\le\sqrt b\,\mathbb E W\).
Thus this gate has null error at most \(64\sqrt b/(\ell M)<1/100\),
and it misses an alternative of mean at most \(3M/4\) with probability
at most the same bound.

If instead the mean exceeds \(3M/4\), let
\(F=\{i\in T:p_i\ge q_i\}\). The score contribution from \(T\setminus F\)
is at most \(\sqrt b\,O_T\le M/8\), and \(q(F)\le O_T\le x\).
Cauchy--Schwarz on \(F\) shows that
\[
  \sum_{i\in T}p_i^2\ge M^2/(4x)=r/4.
\]
The null squared mass is at most \(\sqrt b\,M\le r/8\).

In two new independent batches of size \(\ell\), count matching pairs
with one observation from each batch and with label in \(T\). Reject if
the cross-pair proportion exceeds \(3r/16\). For any law \(p\), put
\(B=\sum_{i\in T}p_i^2\) and \(C=\sum_{i\in T}p_i^3\). The cross count
has mean \(\ell^2 B\). Two indicators for disjoint observation positions
have covariance zero. Indicators that share exactly one position
have joint expectation \(C\), and the identical pair has expectation
\(B\). Counting ordered indicator pairs gives the variance bound
\[
  \operatorname{Var}(\text{cross count})
  \le \ell^2 B+2\ell^3 C
  \le \ell^2 B+2\ell^3 B^{3/2}.
\]
In the last step, every \(p_i\le\sqrt{B}\) on \(T\).
By Chebyshev, the cross-pair gate has null error at most
\(32/(\ell^2 r)+64/(\ell\sqrt r)\), and miss probability
at most \(16/(\ell^2 B)+32/(\ell\sqrt B)\) when \(B\ge r/4\).
Both are below \(1/100\) for \(\ell\sqrt r\ge8192\).
The combined test on \(3\ell\) draws therefore has null error at most
\(1/50\) and miss probability at most \(1/100\).
This is the standard collision moment argument on a selected
subprobability, with a separate weighted score to enforce power
for arbitrary nonuniform alternatives \cite{paninski2008}.

\subsection{Close to the maximum distance}

Let \(\delta=1-e>0\) and
\[
 J=\sum_i q_i\min(1,\sqrt{q_i/\delta}).
\]
When \(J\ge16\delta\), split all labels at cap \(b=\delta\). By
\eqref{eq:upper-clone-moment} the refined moment is at least
\(\sqrt{\delta}\,J/\sqrt 2\ge8\delta\sqrt\delta\).
A strict alternative has full overlap below \(\delta\).
The capped overlap test with \(x=\delta\) uses
\(3\lceil16384/J\rceil\) observations or fewer, including the randomized
refinement. This proves the upper bound \(O(1/J)\) in this branch.

\subsection{The regular regime}

Use the shortest descending prefix \(H\) with \(q(H)\ge\delta\) and write
\(A=q(H)=1-w\). Choose \(i_*\) to be its largest-probability label and put
\(O=H\setminus\{i_*\}\). The head term in original probability units is
\[
 P(n)=\min_{\substack{R\ \mathrm{prefix}\\\mathrm{of}\ O}}
  \left\{\frac{\big(\sum_{i\in R}q_i^{2/3}\big)^{3/4}}{\sqrt n}
              +q(O\setminus R)\right\},\qquad
 D(n)=\sqrt{wA/n}+\frac1n+P(n).
\]
For the normalized conditional head \(h=q_H/A\), this is
\(P(n)=A P_h(nA)\) by direct scaling. No conditional observations are
required for the test below.

First test the tail mass using \(n\) full observations. For the tail count,
use the threshold \(512L\) about \(nw\), where \(L=\sqrt{nwA}+1\).
Its null variance is \(nwA\), giving null error less than \(1/100\).
At alternative mass \(s\), the variance is at most \(n(wA+|s-w|)\),
since \(s(1-s)-w(1-w)=(s-w)(1-s-w)\).
If \(n|s-w|>1024L\), accepting requires a deviation from the alternative
mean of at least \(n|s-w|/2\). Writing \(d=|s-w|\), the resulting miss bound is
\(4(nwA+nd)/(nd)^2\le4/1024^2+4/1024<1/100\).
In particular the same test detects
\(p(S)>4096(q(S)+1/n)\) for an arbitrary prescribed set \(S\), on
replacing \(T\) by \(S\). To see the implication, put \(u=nq(S)\) and
use \(\sqrt{u}+1\le u+2\).

For the remaining head shape, choose a minimizing prefix \(R\) and set
\[
 W_R=\sum_{i\in R}q_i^{2/3},\qquad B_R=W_R^{3/4}/\sqrt n,
 \qquad v=\sqrt{W_R}/n.
\]
If \(W_R=0\), take the retained set \(K\) empty. Otherwise let
\(K=\{i\in R:n^2v q_i^{2/3}\ge1\}\). The extra mass omitted by this
trimming is at most \(W_R/(n\sqrt v)=B_R\).
Take weights \(u_i=q_i^{-2/3}\) on \(K\), zero elsewhere, and define
\[
 Z=\sum_i u_i (N_i-nq_i)(N'_i-nq_i),\qquad
 V=\sum_i u_i(p_i-q_i)^2,
\]
from independent \(n\)-sample histograms \(N,N'\) of the full law.
This is the weighted centered quadratic method used in identity
testing \cite{valiantvaliant2017,chhorcarpentier2022}, here in a fixed
two-batch form. Its expectation is \(n^2 V\). Writing
\(\Sigma=\operatorname{diag}(p)-pp^\top\) for the single-observation
covariance, the centered bilinear part has variance
\(n^2\sum_{i,j}u_i u_j\Sigma_{ij}^2\le2n^2\sum_i u_i^2p_i^2\).
Indeed the diagonal term is at most \(\sum_i u_i^2p_i^2\) and the
off-diagonal term is at most \((\sum_i u_i p_i^2)^2\), which is bounded
by the same quantity. The linear parts are uncorrelated with the
bilinear part and have combined variance at most
\(2n^3\sum_i u_i^2p_i(p_i-q_i)^2\). In particular
\[
 \operatorname{Var}(Z)
 \le 3n^2\sum_i u_i^2p_i^2+
       6n^3\sum_i u_i^2p_i(p_i-q_i)^2.
\]
The cap on \(K\) gives \(\max u_i\le n^2v\), while
\(\sum_i u_i^2q_i^2\le W_R=n^2v^2\) and \(\max u_iq_i\le nv\).
It follows that
\[
 \sum_i u_i^2p_i^2\le2W_R+2n^2vV,\qquad
 \max_i u_i p_i\le n v+n\sqrt{vV}.
\]
Rejecting for \(Z>32n^2v\) gives null error at most \(3/32^2\).
If \(\sum_{i\in K}|p_i-q_i|>10000B_R\), Cauchy--Schwarz and
\(B_R^2=W_Rv\) give \(V>10^8v\). The threshold is below half the
mean. The ratio of variance to squared mean is bounded by
\[
 6(v/V)^2+12(v/V)+6\sqrt{v/V}.
\]
Four times this bound is below \(1/100\), giving the stated miss
probability by Chebyshev.

Apply the mass gate on another \(n\) observations to \(O\setminus K\).
Its null mass is at most \(B_R+q(O\setminus R)=P(n)\).
Unless one of the two shape guarantees applies, we have
\[
 \sum_{i\in O}|p_i-q_i|
 \le 14097P(n)+4096/n .
\]
For \(W_R=0\) the retained gate is unnecessary and the same bound holds.
Using \(|p(T)-w|+\sum_{i\in O}|p_i-q_i|\) to bound \(L_H(p,q)\)
follows by eliminating the single coordinate \(a\) with normalization.
Consequently any strict alternative outside the head rejection guarantees
has \(L_H(p,q)\le G\), where
\[
 G=1024\left(\sqrt{wA/n}+1/n\right)+14097P(n)+4096/n
 \le 14097D(n).
\]

Write \(Y(n)=\Delta+X(n)\), and suppose \(4000000D(n)\le Y(n)\).
Then \(256G\le Y(n)\). If \(X(n)\le\Delta\), this implies \(G\le\Delta\)
and every strict alternative is covered by a head guarantee.
Otherwise put \(x=X(n)/256\). Then \(G\le\Delta+x\), so by
\eqref{eq:upper-head-tail} all remaining strict alternatives have
\(O_T(p,q)<x\). Also \(nx\ge1\), since \(D(n)\ge1/n\) and
\(Y(n)<2X(n)\).

On the tail split at cap \(b=1/(n^2x)\), keeping the other labels
unchanged. The defining root gives
\[
 \sum_{i\in T}\min(q_i/x,nq_i^{3/2}/\sqrt x)\ge16
\]
by comparing each term at \(x=X(n)/256\) to its value at \(X(n)\).
Equation \eqref{eq:upper-clone-moment} then gives a refined moment
\(M\ge8\sqrt x/n=8x\sqrt b\).
Since \(b\le x\), the capped overlap test applies and needs at most
\(3\lceil8192\sqrt x/M\rceil\le3072n\) draws.
The group gate, head gates, and optional tail test use disjoint
batches of full observations. Combining their rejections gives both
errors below \(1/3\) with at most \(3076n\) observations. All inequalities
use the original \(e\), including equality at the population head gate.

\section{Lower bounds from shared pilot lists}
\label{sec:pilot-lower}

We use the half-\(\ell_1\) convention for total variation. Two laws
on the space of ordered records with TV less than \(1/3\) cannot have
rejection probabilities at most \(1/3\) and at least \(2/3\), respectively.
This remains true when the second law mixes product laws of strictly
far alternatives, since the power bound then holds under the mixture.

\paragraph{A lower bound near the maximum distance.}
Put \(\delta=1-e>0\) and
\[
 J_\delta(q)=\sum_i q_i\min\{1,\sqrt{q_i/\delta}\}.
\]
We prove that a successful positive budget \(n\) satisfies
\(nJ_\delta(q)>1/64\). Suppose instead \(nJ_\delta(q)\leq 1/64\).
Remove the atoms with \(q_i>\delta/16\). Their total mass \(R\) is
at most \(4J_\delta(q)\), so \(R\leq 1/16\). Write \(q^\circ\) for
the remaining reference normalized by \(1-R\).
The original and normalized ordered product laws differ in TV by at
most \(1-(1-R)^n\leq nR\leq 1/16\). This follows by conditioning the
original record on avoiding all removed labels.

Partition the retained positive atoms into occupied dyadic bins
\((u_j/2,u_j]\) with \(u_j=(\delta/16)2^{-j}\). Let \(b_j\) be the
bin mass under \(q^\circ\), and put \(M=\sum_j b_j\sqrt{u_j}\).
For a retained atom \(q_i\) in bin \(j\), we have
\(\sqrt{u_j}\leq 2\sqrt{q_i}\). Thus
\[
 \sqrt{\delta}\,M
 \leq \frac{2\delta}{1-R}
          \sum_{i:\,q_i\leq \delta/16}q_i\sqrt{q_i/\delta}
 \leq 4\delta J_\delta(q).
\]
We also have \(\sum_j u_j\leq \delta/8\).
For each bin draw independent pilot labels from its conditional
reference, with list length
\[
 k_j=\left\lceil n\sqrt{\delta}\,b_j/\sqrt{u_j}\right\rceil.
\]
A realized alternative allocates mass \(b_j\) to the uniform
distribution on the list slots in bin \(j\), counting repeated
labels with their multiplicities. Each realized law is normalized.
The reference mass of its support is at most
\[
 \sum_j k_j u_j
 \leq n\sqrt{\delta}\,M+\sum_j u_j
 \leq \delta/16+\delta/8<\delta.
\]
Its overlap with \(q\) is therefore strictly less than \(\delta\).

In the ordered predictive law, each coordinate independently
chooses its bin and a uniform slot from the shared list. Couple
this law to \((q^\circ)^{\otimes n}\) by using the corresponding
pilot label on the first use of each slot and fresh independent
conditional draws for repeat uses in the reference record. A
difference can occur only when a pair uses the same bin and slot.
The resulting TV bound is
\[
 \binom n2\sum_j\frac{b_j^2}{k_j}
 \leq \frac{n}{\sqrt{\delta}}M
 \leq 4nJ_\delta(q)\leq 1/16.
\]
Together the two TV costs are at most \(1/8\), which yields the
contradiction. If the strict alternative is nonempty, zero draws
cannot give the two errors, so the bound also provides a positive
sample-scale lower bound in cases where \(1/(64J_\delta(q))<1\).

\paragraph{The regular pilot construction.}
Let \(H\) be the chosen head, \(T\) its tail, \(w=q(T)\), and
\(\Delta=e-w\). For any full law \(p\), define
\[
 L(p)=\sum_{i\in H}(q_i-p_i)_+,\qquad
 O(p)=\sum_{i\in T}\min(q_i,p_i).
\]
The exact identity is \(\operatorname{TV}(p,q)=w+L(p)-O(p)\).
Keep the head labels and replace \(T\) by a single marker, writing
\(\bar q\) for this reduced reference.

Suppose \(w>0\), \(n\geq1\), and \(x>0\) obeys
\[
 \rho=K_T(n,x)
  :=\sum_{i\in T}\min\{q_i/x,\ n q_i^{3/2}/\sqrt{x}\}
  \leq 1/64,\qquad nx\geq 8.
\]
Freeze the tail atoms with \(q_i\geq \tau=1/(n^2x)\).
The two terms in the minimum cross at \(\tau\), giving
\[
 \sum_{i\in T:\,q_i\geq\tau}q_i\leq \rho x,\qquad
 n\sqrt{x}\sum_{i\in T:\,q_i<\tau}q_i^{3/2}\leq \rho x.
\]
Put the other positive atoms in occupied dyadic bins
\((u_j/2,u_j]\) with \(u_j=\tau 2^{-j}\). Denote their
\emph{unconditional} masses under \(q\) by \(b_j\), and set
\(M=\sum_j b_j\sqrt{u_j}\). Then
\[
 n\sqrt{x}\,M\leq 2\rho x,\qquad
 \sum_j u_j\leq 2\tau\leq x/32.
\]
For each bin draw an independent list of
\[
 k_j=\left\lceil n\sqrt{x}\,b_j/\sqrt{u_j}\right\rceil
\]
labels from the conditional reference in that bin.

The channel takes an ordered reduced record. It first draws the lists,
independently of the input record. It keeps every head label. For
each marker it independently chooses a frozen label \(i\) with
weight \(q_i/w\), or a bin \(j\) with weight \(b_j/w\) followed by
a uniform slot in that list. This prescription defines one common
Markov channel on ordered records. Conditional on the lists, an
IID reduced input yields an IID full output. For any reduced input
law and every realized list collection, the full tail overlap is
at most the reference mass of the frozen labels and the occupied
slots. This is bounded by
\[
 \rho x+\sum_j k_j u_j
 \leq 3\rho x+x/32\leq x.
\]
In particular, this bound holds for every possible probability of
the marker in a reduced alternative.

Under \(\bar q^{\otimes n}\), each light bin is selected with its
unconditional probability \(b_j\). The same first-use coupling
works on ordered records, leaving the head and frozen draws
fresh. The predictive TV cost of the fill under this reduced
null is bounded by
\[
 \binom n2\sum_j\frac{b_j^2}{k_j}
  \leq \frac{nM}{\sqrt{x}}\leq 2\rho .
\]
Let a prior on reduced laws have predictive TV at most \(1/8\)
from \(\bar q^{\otimes n}\), and suppose each positive-prior reduced
law has head deficit strictly larger than \(\Delta+x\). The
full filled components all have strict TV greater than \(e\).
By data processing for the common channel and the preceding
null comparison, their predictive TV from \(q^{\otimes n}\) is
at most \(1/8+2\rho<1/3\). No full test with both errors at most
\(1/3\) succeeds.

If \(X(n)\) is the positive root of \(K_T(n,X)=1\), take
\[
 x=4096\max\{X(n),1/n\}.
\]
This meets the two scalar conditions. Each summand shrinks by at
least a factor of \(64\) when the second argument is increased by
a factor of \(4096\). Labels of zero null mass have zero pilot mass. If \(w=0\),
a reduced head prior whose marker remains zero transfers by
keeping its head draws, at zero additional predictive cost. Full
tail overlap is then zero, including when the declared tail is
empty.

\paragraph{The marker fluctuation.}
Write \(A=1-w\). For \(w>0\) and \(nA\geq1\), let
\[
 d=\frac18\sqrt{wA/n}.
\]
Since \(w\leq1\leq nA\), we have \(d\leq A/8\).
Consider the reduced law whose marker has mass \(w+d\) and
whose head probabilities are \(q_i(1-d/A)\). This is a probability
law, with head deficit exactly \(d\). Its one-draw chi-square
divergence with respect to \(\bar q\) is
\[
 d^2/w+d^2/A=d^2/(wA)=1/(64n).
\]
Its ordered product TV from the reduced null is bounded by
\[
 \frac12\sqrt{(1+1/(64n))^n-1}
 \leq \frac{1}{2\sqrt{63}}<\frac1{12}.
\]
The middle inequality follows from the binomial expansion and
\(\binom nk\leq n^k\), followed by the geometric sum.
This singleton reduced prior meets the predictive condition
above. If a full test succeeds, it follows that \(d\leq\Delta+x\),
or
\[
 \sqrt{w(1-w)/n}\leq 8(\Delta+x).
\]
When \(w=0\) the left-hand side vanishes and the same necessity
holds for \(\Delta\geq0\).

\section{Signed perturbations on the retained head}
\label{sec:signed-head-lower}

The retained \(\ell_{2/3}\) term, deletion of a largest atom, and
clipped \(2/3\)-power shifts appear in Valiant and Valiant's
instance-optimal testing analysis \cite{valiantvaliant2017}.
Chhor and Carpentier formulate the corresponding fixed-sample
separation radius and use signed priors in their bulk lower bound
\cite[Corollary~1 and Section~5.1]{chhorcarpentier2022}. We use
capped, paired shifts and keep the original head probabilities. This
preserves the total head mass and lets the same prior be composed with
the independent tail construction.

List the positive head probabilities as \(q_1\ge\cdots\ge q_d>0\).
They need not sum to one. Set \(w=1-\sum_{i=1}^{d}q_i\) and
\(\Delta=e-w\). On the reduced alphabet, write \(q^\sharp\) for the
law with these \(d\) probabilities and an additional tail marker of
mass \(w\). For \(n\ge1\), use the head profile
\begin{equation}
 P(n)=\min_{1\le j\le d}
 \left\{
  \frac{\left(\sum_{i=2}^{j}q_i^{2/3}\right)^{3/4}}{\sqrt n}
  +\sum_{i=j+1}^{d}q_i
 \right\}.
 \label{eq:signed-head-profile}
\end{equation}
If the normalized head law is \(h_i=q_i/A\), where \(A=1-w>0\),
then this profile equals \(A P_h(nA)\) with \(P_h\) given by the
same formula on \(h\). Both terms scale by \(A\) under this
substitution. We will use \(P(n)\) directly.

\paragraph{A capped shift with a quartic budget.}
For \(i=2,\ldots,d\), choose
\[
 z_i=\min(q_i,\lambda q_i^{2/3}),
 \qquad
 \sum_{i=2}^{d}\frac{z_i^4}{q_i^2}\le\frac1{n^2}.
\]
If \(\sum_{i=2}^{d}q_i^2\le n^{-2}\), take \(\lambda\) large
enough that every \(z_i=q_i\). Otherwise the budget is attained with
equality for some \(\lambda>0\), by continuity between
\(\lambda=0\) and \(\lambda=\max_{i\ge2}q_i^{1/3}\). In the
singleton case all sums are empty and the saturated choice applies.
We extend the same definition of \(z_i\) to \(i=1\) for ordering
purposes and write \(S=\sum_{i=2}^{d}z_i\).

For any cutoff \(j\), H\"older's inequality and the quartic budget
give
\[
 \sum_{i=2}^{j} z_i
 \le \left(\sum_{i=2}^{j}z_i^4/q_i^2\right)^{1/4}
       \left(\sum_{i=2}^{j}q_i^{2/3}\right)^{3/4}
 \le n^{-1/2}\left(\sum_{i=2}^{j}q_i^{2/3}\right)^{3/4}.
\]
The remaining shift mass is at most the remaining reference mass.
Thus \(S\le P(n)\). In the all-saturated case the cutoff \(j=1\)
also gives \(P(n)\le S\). Otherwise put
\(I=\{i\ge2:q_i>\lambda^3\}\) and
\(I^c=\{2,\ldots,d\}\setminus I\). The retained indices \(I\)
form a prefix. Write
\[
 W=\sum_{i\in I}q_i^{2/3},\qquad
 U=\sum_{i\in I^c}q_i,\qquad R=\sum_{i\in I^c}q_i^2.
\]
Here \(S=\lambda W+U\), \(n^{-2}=\lambda^4W+R\), and
\(R\le\lambda^3 U\). The fourth power of the bulk term at this
cutoff satisfies
\[
 \frac{W^3}{n^2}
 \le (\lambda W)^4+(\lambda W)^3U
 \le (\lambda W+U)^4.
\]
Consequently \(S\le P(n)\le2S\). Atoms tied at
\(\lambda^3\) may stay outside the prefix.

\paragraph{The paired prior.}
Use the disjoint pairs \((1,2),(3,4),\ldots\). For each pair \(j\)
draw an independent uniform sign \(\sigma_j\in\{-1,1\}\) and add
\(\sigma_j b_j\) to coordinate \(2j-1\) and subtract it from
coordinate \(2j\), where \(b_j=z_{2j}/8\). Write \(p_\sigma^\sharp\)
for the resulting reduced law. Since \(b_j\le q_{2j}\le q_{2j-1}\),
the shifts give a nonnegative normalized law. They leave the marker
and any unpaired label unchanged. In particular the construction
permits a marker of mass zero. The function
\(t\mapsto\min(t,\lambda t^{2/3})\) is nondecreasing, so the \(z_i\)
are sorted. Each odd index \(i\ge3\), including a possible unpaired
last index, has \(z_i\le z_{i-1}\). For every sign vector the
unconditional head deficit is therefore
\begin{equation}
 L_H(p_\sigma^\sharp,q^\sharp)=\sum_j b_j
   \ge S/16\ge P(n)/32 .
 \label{eq:signed-head-deficit}
\end{equation}

Let \(M=\mathbb E_\sigma (p_\sigma^\sharp)^{\otimes n}\) and
\(Q=(q^\sharp)^{\otimes n}\) on ordered \(n\)-tuples. The
perturbations have no mass outside the null support, including at a
zero marker. With independent sign vectors \(\sigma,\tau\), product
factorization of the likelihood inner product gives
\begin{align*}
 1+\chi^2(M,Q)
 &=\mathbb E_{\sigma,\tau}
   \left(\sum_{a:q^\sharp_a>0}
      \frac{(p_\sigma^\sharp)_a(p_\tau^\sharp)_a}{q^\sharp_a}
   \right)^n  \\
 &=\mathbb E_{\sigma,\tau}
   \left(1+\sum_j c_j\sigma_j\tau_j\right)^n,
 & c_j&=b_j^2\left(\frac1{q_{2j-1}}+\frac1{q_{2j}}\right).
\end{align*}
The linear shift terms cancel within each pair, and the unchanged
labels contribute their null masses to the constant term. The
quantity raised to the power \(n\) is nonnegative. The inequality
\(1+u\le e^u\) and independence of the sign products now give
\[
 1+\chi^2(M,Q)
 \le \prod_j\cosh(nc_j)
 \le \exp\!\left(\frac{n^2}{2}\sum_j c_j^2\right)
 \le \exp(1/2048).
\]
The middle step follows from the power series and
\((2k)!\ge2^k k!\). The last follows from
\(c_j\le 2b_j^2/q_{2j}\) and the quartic budget. Thus
\(\operatorname{TV}(M,Q)\le\frac12\sqrt{e^{1/2048}-1}<1/8\)
by Cauchy--Schwarz. For the numerical inequality one can use
\(e^t\le(1-t)^{-1}\) at \(t=1/2048\).

\paragraph{Consequences at the exact radius.}
Suppose the tail construction gives a prior of tail laws of mass
\(w\) whose reference overlap is at most \(x\) for every draw of the
prior. Suppose also that when the head probabilities are unchanged,
its ordered \(n\)-sample mixture is within \(1/8\) of the full
null. When \(x\ge w\), the unchanged tail is such a prior. In
particular the argument covers \(w=0\), where the marker is never
observed. For positive \(w\), choose the tail prior independently
of the head signs and replace each marker by an independent draw
from the chosen conditional tail law, given that law. This is a
single common randomized channel on ordered \(n\)-tuples. It
contracts total variation, so the full mixture of paired alternatives
is within \(1/8+1/8=1/4\) of the null.

If \(P(n)>32(\Delta+x)\), equation~\eqref{eq:signed-head-deficit}
gives \(L_H>\Delta+x\) for every paired alternative. Marker
replacement leaves this deficit unchanged and gives
\(O_T=\sum_{i\in T}\min(p_i,q_i)\le x\). Hence every lifted law
has distance \(w+L_H-O_T>e\). A test with each of its two errors
at most \(1/3\) would then have a rejection-probability gap of at
least \(1/3\) between the null and this mixture. That is impossible
at mixture TV at most \(1/4\). Thus a successful \(n\)-draw test
necessarily satisfies
\begin{equation}
 P(n)\le32(\Delta+x)
 \label{eq:signed-head-necessity}
\end{equation}
whenever the stated tail-prior bounds hold. The radius in this
comparison is the original strict threshold \(e\).


\section{Comparing sample budgets at the same distance}\label{sec:horizontal}
This section completes the finite proof. The near-one tests and the small-$J$ lower construction give $mJ>1/64$ for every successful budget $m$ on a nonempty instance, and a test with $3\lceil16384/J\rceil$ samples whenever $J\ge16\delta$. Rounding gives $\Phi\le65m$ and an upper bound of $49152\Phi$ on this branch.

In the other branch we can use the slightly weaker assumption $J\le32\delta$. Suppose $m$ observations suffice. A strictly far witness $p$ can be interpolated with $q$ to any prescribed distance strictly between $e$ and $\TV(p,q)$. The product total variation is at most $m$ times the one-observation distance. Letting the latter decrease to $e$ gives $me\ge1/3$. The preceding bound on $mJ$ gives $m\delta>1/2048$. Put $n=2^{43}m$. A test can discard its extra observations, so $n$ is a successful budget as well.

Put $x=4096\max\{X(n),1/n\}$. The preceding signed-head and marker bounds give
\[
       P(n)\le32(\Delta+x),\qquad
       \sqrt{wA/n}\le8(\Delta+x).
\]
The marker estimate requires $nA\ge1$, which follows from $A\ge\delta$. If $nY(n)<1$, then $X(n)<1/n$ and $n\Delta<1$. The marker bound gives $nwA\le32776^2$. If $w\le1/2$, it follows that $ne=nw+n\Delta\le2(32776)^2+1$. If $w>1/2$, then $n\delta\le nA\le2(32776)^2$. Either conclusion contradicts $ne,n\delta\ge2^{32}$. Thus $nY(n)\ge1$ and $x\le4096Y(n)$. In particular
\begin{equation}\label{eq:weak}
             D(n)\le163881\,Y(n).
\end{equation}

\begin{lemma}\label{lem:dilate}
For each $K,c>0$ there is a finite $M=M(K,c)$, independent of the law, such that $D(n)\le K Y(n)$ with $n\ge1$ implies $D(n')\le cY(n')$ at a power of two $n'$ in $[n,2Mn]$.
\end{lemma}
\begin{proof}
Both $D$ and $-Y$ are nonincreasing. Since $D(n)\ge1/n$, a weak crossing implies $nY(n)\ge G:=1/K$. Choose $M_1\ge1$ with $2K\le c M_1^2$. Choose $\gamma>0$ with $2\gamma\le G^2$ and $M_1^4\gamma\le1$. Finally take $t\ge1$ with $2(K+1/\gamma)\le ct$.

For any omitted head suffix of mass $u$, the sum $W$ of its $2/3$ powers is at most $u/a^{1/3}$. Consequently $W^{3/4}/\sqrt n\le u+1/(an^2)$, by considering $an^2u\ge1$ and its complement. At a minimizing head cutoff in \eqref{eq:P}, append this omitted suffix and use $(v+W)^{3/4}\le v^{3/4}+W^{3/4}$. The full head, evaluated at $t^2n$, gives
\[
        D(t^2n)\le\frac2t\big(D(n)+1/(an^2)\big).
\]
If $n^2aY(n)\ge\gamma$, this gives $D(t^2n)\le cY(t^2n)$.

For the other case write $Y=Y(n)$. If $Y>2X(n)$ then $Y<2\Delta\le2a$, whence $(nY)^2<2\gamma\le G^2$, a contradiction. Thus $X(n)\ge Y/2>0$. For every tail atom, $M_1^4 n^2q_iX(n)\le M_1^4 n^2aY<1$. At $M_1n$ all terms of \eqref{eq:K} at $x'=M_1^2 X(n)$ are on their unsaturated branch, as they were at $(n,X(n))$. Their sum is still one, so $X(M_1n)=M_1^2X(n)$. Therefore $D(M_1n)\le D(n)\le KY\le cY(M_1n)$. Increase the successful real budget to the next power of two. Monotonicity preserves the inequality, and the increase is at most a factor of two. Take $M=\max(M_1,t^2)$.
\end{proof}

Apply Lemma~\ref{lem:dilate} to \eqref{eq:weak} with $K=163881$ and $c=c_*$. The first dyadic crossing is at most $2^{44}M(163881,c_*)m$. Conversely the regular test at a strong crossing uses at most $3076n$ original observations. These bounds prove the finite statistical assertion of Theorem~\ref{thm:main}.

\section{An explicit finite-arithmetic rate}
\label{sec:closed-rate}

This section gives a root-free expression for the rate on a finite
rational input and describes its bit complexity. Let
\(\lfloor z\rfloor_2\) be the largest integral power of two at most
\(z>0\), and let \(\lfloor0\rfloor_2=0\). The empty-alternative branch
returns zero when \(e\ge1-\min_iq_i\). The minimum is taken on the
declared alphabet, including zero-probability labels.

Assume the alternative is nonempty, set \(\delta=1-e\), and compute
\[
  \widehat J=\sum_i q_i\min\{1,\lfloor\sqrt{q_i/\delta}\rfloor_2\}.
\]
Then \(J/2\le\widehat J\le J\). If \(\widehat J\ge16\delta\), return
\(\lceil1/\widehat J\rceil\). This is the near-maximum-distance
rate, since in this branch \(J\ge16\delta\). In the other branch
\(J<32\delta\), which is sufficient for all the regular-regime
baseline comparisons.

Use the shortest decreasing-mass prefix \(H\) with \(q(H)\ge\delta\).
Its complement \(T\) has mass \(w\). Put \(A=1-w\), \(\Delta=e-w\),
and \(O=H\setminus\{i_*\}\), with \(i_*\) a largest head label.
The head cuts \(R\) are prefixes of \(O\), and the tail cuts \(S\) are
prefixes of \(T\), each in decreasing reference mass. It suffices
to use endpoints of equal-probability runs, including the empty
and full cut in each list. For each \(R\) set
\[
 \begin{split}
 c_{R,0}&=q(O\setminus R),\\
 c_{R,1}&=\max\left\{\lfloor\sqrt{wA}\rfloor_2,\
       \left\lfloor
          \left(\sum_{i\in R}\lfloor q_i^{2/3}\rfloor_2\right)^{3/4}
       \right\rfloor_2\right\},\\
 c_{R,2}&=1.
 \end{split}
\]
For each \(S\) set
\[
 b_S=\max\{\Delta,q(S)\},\qquad
 h_S=\left(\sum_{i\in T\setminus S}q_i\lfloor\sqrt{q_i}\rfloor_2\right)^2 .
\]
A tail consisting solely of null-zero atoms contributes only the
empty cut, with \(b_S=\Delta\) and \(h_S=0\).

For \(c,b\ge0\) and a nonnegative integer \(j\), define
\[
 E(c,b,j)=
 \begin{cases}
 0,&c\le b,\\
 +\infty,&c>b\ \text{and}\ (b=0\ \text{or}\ j=0),\\
 \displaystyle\left\lceil\frac{\lceil\log_2(c/b)\rceil}{j}\right\rceil,
   &\text{otherwise}.
 \end{cases}
\]
The resulting rate in the regular regime is
\begin{equation}
 4^{\,\displaystyle\min_R\,
       \max_{\substack{S\\ a\in\{0,1,2\}}}
       \min\{E(c_{R,a},b_S,a),\,E(c_{R,a},h_S,a+4)\}} .
 \label{eq:closed-rate}
\end{equation}
The use of \(a\) as a numerical index inside \eqref{eq:closed-rate}
is local to that display.

\subsection{Comparison with the statistical crossing}

For any \(n\ge1\) define
\[
 \widehat D(n)=\min_R\max_{a=0,1,2}c_{R,a}n^{-a/2},
 \qquad
 \widehat Y(n)=\min_S\max\{b_S,h_Sn^2\}.
\]
The original head cost for each \(R\) is the sum of four terms:
the omitted mass, the group standard deviation divided by \(\sqrt n\),
the retained bulk divided by \(\sqrt n\), and \(1/n\).
The group numerator loses at most a factor 2 in rounding. The
head numerator loses at most \(2^{7/4}<4\), including the inner
rounding and the final rounding. Thus at every cut the rounded
maximum lies between one eighth of the original cost and the
original cost. Within an equal-probability run, the original head
cost is affine plus a positive multiple of a \(3/4\)-power of an
affine function of the number of retained labels. Concavity places
its minimum at an endpoint. Consequently
\begin{equation}
  D(n)/8\le\widehat D(n)\le D(n).
  \label{eq:closed-head-comparison}
\end{equation}
The rounded maximum itself may prefer a cut within a tied run.
The same constant-factor estimate holds with endpoint compression.
An expanded scan of all tied labels can return a different integer.

For the tail, let \(z=\sqrt{X(n)}\), and put
\(Q_S=q(S)\), \(H_S=\sum_{i\in T\setminus S}q_i^{3/2}\).
If \(z>0\), its defining equation gives
\[
 z^2=\sum_{i\in T}\min(q_i,nq_i^{3/2}z)
       \le Q_S+nH_Sz
       \qquad\text{for every }S.
\]
If \(B_S=\max\{Q_S,(nH_S)^2\}\), this implies
\(z^2\le B_S+\sqrt{B_S}\,z\), hence \(z^2\le4B_S\).
Conversely, equality holds in the previous display for the prefix
\(\{i:nz\sqrt{q_i}\ge1\}\), which then has \(B_S\le z^2\).
This prefix ends at an equal-run boundary, including when some
atoms are exactly at the saturation threshold. Rounding the square
root in \(H_S\) costs at most a factor 2, so that
\[
 \widehat X(n):=\min_S\max\{Q_S,h_Sn^2\}\in[X(n)/16,X(n)].
\]
If the tail mass is zero, all these tail moments and \(X\) vanish.
Taking the maximum with \(\Delta\) gives
\begin{equation}
  Y(n)/17\le\widehat Y(n)=\max\{\Delta,\widehat X(n)\}\le Y(n),
  \qquad Y(n)=\Delta+X(n).
  \label{eq:closed-tail-comparison}
\end{equation}

For \(n=4^k\), the condition \(\widehat D(n)\le\widehat Y(n)\)
means that some \(R\) satisfies, for every \(S\) and \(a=0,1,2\),
at least one of
\[
 c_{R,a}\le b_S2^{ak},
 \qquad c_{R,a}\le h_S2^{(a+4)k}.
\]
Each inequality is an upward interval of nonnegative integer
exponents. Its least exponent is \(E\). The first exponent in
the union of two such intervals is the smaller of the two
\(E\)'s, and their intersection over the requirements has first
exponent equal to the maximum. Minimizing over \(R\) gives
\eqref{eq:closed-rate}. At the full head cut the constant
omitted-mass term is zero, so a finite exponent always exists:
for every tail cut at least one budget is positive when \(e>0\).

A successful closed comparison enforces the baseline scales
even without an explicit lower search floor. The full tail cut
forces \(1/n\le\max(\Delta,w)\le e\).
For the other scale, choose \(S=\{i\in T:q_i\ge\delta\}\).
In the regular branch \(q(S)\le J<32\delta\), and
\(\Delta<\min_{i\in H}q_i\le\max_iq_i\le\max(J,\delta)\).
The remaining moment is at most \(\sqrt\delta J\). The \(1/n\)
term thus forces
\[
 1/n\le\max\{32\delta,1024\delta^3n^2\},
\]
and hence \(n\delta\ge1/32\).

By \eqref{eq:closed-head-comparison}--\eqref{eq:closed-tail-comparison},
passing the closed comparison implies \(D(n)\le8Y(n)\).
Conversely, \(D(n)\le Y(n)/17\) suffices to pass. The horizontal
comparison and the statistical bounds therefore identify
\eqref{eq:closed-rate} with the sample complexity up to universal
constants at the same radius. A fixed dilation covers the baseline
floor, and rounding a clock up to the next power of four costs
at most a factor four.

\subsection{Exact arithmetic and scope}

A finite input may specify distinct rational atom weights together
with their binary multiplicities. Sort the weights, and locate the
shortest head by one rational ceiling if it ends within a run.
Delete one largest atom for the head cost. The remaining head
and tail scans use run endpoints. The displayed sums are computed
using multiplicities and do not enumerate the labels in a run.

A downward rounding \(\lfloor x^{u/v}\rfloor_2\) for the fixed
rational powers used here is computed by comparing the numerator and
denominator of \(x^u\) with integer shifts. Similarly, the ceiling binary
logarithm of a positive rational ratio is its numerator-denominator
bit-length difference adjusted by one exact comparison. The
definition of \(E\) then uses integer ceiling division. Common
denominators have bit length bounded by the sum of the input
denominator bit lengths, and each fixed-power and shift step
preserves polynomial bit length. The exponent and the binary
output \(4^k\) have polynomial bit length. The number of head-tail
cut pairs is quadratic in the number of input runs. Thus
\eqref{eq:closed-rate} is evaluable in polynomial bit time without
a root computation, a scan over sample sizes, or expansion of binary
multiplicities. This bit bound concerns evaluation of the rate.
The statistical theorem bounds the number of observations.

For a countable alphabet, the next section gives an exact projection argument
and relates the direct profile to its finite evaluations.

\section{Countable alphabets by finite projections}
\label{sec:countable-projections}

Let the declared alphabet be \(\mathbb N\). For \(m\geq 0\), let
\(\pi_m\) keep the first \(m\) labels distinct and send all remaining
labels to one tail cell. For a law \(p\) write \(p_m=(\pi_m)_*p\).
Every projection retains its tail cell, including when the reference
assigns that cell probability zero.

\paragraph{Exact fixed-budget equivalence.}
For each nonnegative integer \(n\) and a fixed strict radius \(e\),
a test with \(n\) ordered IID observations exists for \(q\) if and only
if a test with that budget exists for every \(q_m\) at the same radius.
The two separate error bounds are preserved.

To see the forward implication, fix \(m\). Within a cell of positive
\(q_m\)-mass, independently lift a projected observation by drawing
from the conditional law of \(q\) in that cell. In a cell of zero
\(q_m\)-mass use a point mass at an arbitrary declared label of that
cell. These conditionals have disjoint supports. Thus a lifted finite
law \(\widetilde r\) satisfies \(\widetilde q_m=q\) and
\[
 \operatorname{TV}(\widetilde r,q)=\operatorname{TV}(r,q_m).
\]
The full test applied to the independently lifted ordered observations
gives a finite randomized test with the required errors.

For the converse, lift a projected test to a rejection function
\(f_m\colon \mathbb N^n\to[0,1]\) by applying \(\pi_m\) coordinatewise.
The cube \([0,1]^{\mathbb N^n}\), in its product topology, is sequentially
compact. Select a subsequence \(m_j\to\infty\) for which \(f_{m_j}\)
converges pointwise to \(f\colon\mathbb N^n\to[0,1]\). This subsequence
depends only on the sequence of tests. Under each fixed law \(p^{\otimes n}\),
bounded convergence yields
\[
 \mathbb E_{p^{\otimes n}} f
   =\lim_{j\to\infty}\mathbb E_{p^{\otimes n}} f_{m_j}.
\]
The null bound holds at every index. For each \(p\) with
\(\operatorname{TV}(p,q)>e\), the singleton contribution
\(\tfrac12\sum_{i<m}|p_i-q_i|\) converges to
\(\operatorname{TV}(p,q)\). The projected distance is at least this
contribution, so it is strictly greater than \(e\) for all sufficiently
large \(m\). Bounded convergence therefore also gives the stated
alternative bound. This argument includes \(n=0\).

\paragraph{Finiteness at positive radius.}
Pooling the tail \(\mathcal T_m\) loses at most \(q(\mathcal T_m)\) in
TV. In fact, the loss equals
\[
 \frac12\left(\sum_{i\in\mathcal T_m}|p_i-q_i|
                  -|p(\mathcal T_m)-q(\mathcal T_m)|\right)
 \leq \min\{p(\mathcal T_m),q(\mathcal T_m)\}.
\]
For \(e>0\), choose \(m\) with \(q(\mathcal T_m)<e/2\). A finite test for
\(q_m\) at radius \(e/2\) then yields a finite-budget test for \(q\) at
radius \(e\). The use of \(e/2\) proves finiteness only and does not enter
the fixed-budget comparison.

Let \(N_q(e)\) denote the least successful budget at the two fixed error
levels. Fixed-budget equivalence and the option to ignore extra
observations give
\[
 N_q(e)=\sup_{m\geq0}N_{q_m}(e).
\]
In particular, if \(\Phi(q,e)\) is the finite-alphabet formula of the
main theorem, then
\[
 \Phi_{\mathbb N}(q,e):=\sup_{m\geq0}\Phi(q_m,e)
\]
characterizes \(N_q(e)\) with the same universal comparison constants.
This defines a distribution-dependent functional on countable laws,
without a finite runtime guarantee for arbitrary representations of
infinite inputs. Any countably infinite declared alphabet can be
enumerated to use this formulation.

\paragraph{Evaluation directly on a countable null.}
The original head and tail expression also gives a profile on the countable law.
Sort the positive atoms and retain all declared null-zero labels in the tail.
If $e=1$ the alternative is empty. If $e<1$, a countably infinite alphabet
has atoms of arbitrarily small mass and admits a strict alternative.
Only finitely many head atoms are needed to reach $1-e$. Include in the
prefix projection every atom of mass at least the last head mass $a>0$,
and take $m$ large enough that the aggregate unexposed mass $r_m<a$.
All later projections have the same head probabilities, $P$, $w$ and
$\Delta$. This also covers finite positive support on an infinite alphabet,
for which a zero-mass tail cell remains available once that support is exposed.

For fixed $n,x>0$, the contribution of a single atom $z$ to $K(n,x)$
is superadditive in $z$ since its ratio to $z$ is nondecreasing.
The same holds for the summands of $J$. Consequently $K_m(n,x)$ and
$J_m$ decrease to the corresponding formulae on the full sequence, with
errors at most $r_m/x$ and $r_m$ respectively. If $w>0$, the continuous
strictly decreasing function $K(n,\cdot)$ has a unique positive unit level.
For every $t>X(n)$, eventually $K_m(n,t)<1$ whereas $K_m(n,X(n))\ge1$.
Thus $X_m(n)\downarrow X(n)$. When $w=0$ all roots vanish. The small-overlap
branch stabilizes even if $J=16\delta$, as $J_m\ge J$. Its rounded reciprocal
stabilizes from below. On the regular branch the first crossing $n_*$ for
the full sequence is finite by the same full-head bound and $Y(1)>0$.
It passes every sufficiently fine projected comparison even at equality.
Every earlier dyadic comparison fails strictly in the limit, so all the
finitely many earlier comparisons eventually fail. Hence $\Phi(q_m,e)$
eventually equals the direct countable formula, while $N_{q_m}(e)$ is a
bounded nondecreasing integer sequence and eventually equals $N_q(e)$.
The finite theorem therefore gives the same comparison constants for the
direct profile. At any positive tail root the saturated prefix in
\eqref{eq:radical} is finite. A countable infimum describes the minimal
tail expression at fixed $n$. No finite representation is assumed for
an arbitrary infinite probability law.

\section{Verification, scope and contributions}
The Lean 4 development constructs randomized tests on ordered independent
observations and proves lower bounds against all such tests. It includes
the empty-alternative criterion, integer sample budgets, both finite
branches, the countable projection equivalence and a uniform sample-size
comparison. The finite head and the finite suffix admit independent label
permutations. On explicitly listed rational inputs a further formal theorem bounds the number of dyadic comparisons by a polynomial in the input bit length and evaluates each comparison using finite prefix radicals. The retained prefix includes its boundary atom. The proof
terms use only Lean's standard propositional extensionality, quotient
soundness and classical choice. A separate Rust kernel, nanoda, checks
the exported declarations without access to the Lean sources or Mathlib.
The reproducibility package pins Lean 4.36.0 and its Mathlib revision.

The finite radical formula and the countable theorem using the supremum
over projections have corresponding Lean declarations. The passage from
the rounded integer arithmetic in Section~\ref{sec:closed-rate} to a
polynomial bit bound, and the eventual equality with the direct countable
profile in Section~\ref{sec:countable-projections}, are given as mathematical
proofs in this paper. The Python implementation has independent exact-arithmetic
checks. The result concerns universal multiplicative constants. Their size
leaves open useful questions about numerical calibration for small samples.

Yoshiaki Otsuka set the research question, the acceptance criteria and
the conditions for public release. OpenAI Codex research agents developed
the mathematical constructions, wrote the manuscript and Lean sources,
and ran the verification. Parallel agents investigated the upper tail
statistic, competing lower constructions, and signed head perturbations.
Separate agent reviews examined fidelity to the source question and
searches for prior equivalent results. The proof does not rely on the
agents' agreement. The public record contains source hashes, the independent
checker receipt, reproducibility instructions and a computation log.
A public IFM dataset index lists a prior generator run and sixteen verification sessions concerning this question. Its public summary gives no question-specific verdict. Our priority audit found no earlier publicly inspectable general solution. No external human peer review had been completed at the time of writing.


\begin{thebibliography}{9}
\bibitem{Canonne} C.~L. Canonne. Open problem: tight characterization of instance-optimal identity testing. COLT, PMLR 247, 5312--5316, 2024.
\href{https://proceedings.mlr.press/v247/canonne24a.html}{Proceedings and Open Question 1}.
\bibitem{valiantvaliant2017} G.~Valiant and P.~Valiant. An automatic inequality prover and instance optimal identity testing. SIAM J. Comput. 46(1), 429--455, 2017. \href{https://doi.org/10.1137/151002526}{DOI: 10.1137/151002526}. Conference version, FOCS 2014.
\bibitem{BCG} E.~Blais, C.~L. Canonne and T.~Gur. Distribution testing lower bounds via reductions from communication complexity. CCC, LIPIcs 79, 28:1--28:40, 2017.
\href{https://doi.org/10.4230/LIPIcs.CCC.2017.28}{DOI: 10.4230/LIPIcs.CCC.2017.28}.
\bibitem{chhorcarpentier2022} J.~Chhor and A.~Carpentier. Sharp local minimax rates for goodness-of-fit testing in multivariate binomial and Poisson families and in multinomials. Math. Stat. Learn. 5(1/2), 1--54, 2022. \href{https://doi.org/10.4171/MSL/32}{DOI: 10.4171/MSL/32}.
\bibitem{paninski2008} L.~Paninski. A coincidence-based test for uniformity given very sparsely sampled discrete data. IEEE Trans. Inform. Theory 54(10), 4750--4755, 2008. \href{https://doi.org/10.1109/TIT.2008.928987}{DOI: 10.1109/TIT.2008.928987}.
\bibitem{Zipf} P.~Dewasurendra. Closing the harmonic gap: exact identity testing for Zipfian distributions. Author-posted manuscript, 2026.
\href{https://andotheror.github.io/p/closing-the-harmonic-gap-exact-identity-testing-for-zipfian-distributi.html}{Author's manuscript}.
\end{thebibliography}
\end{document}
